\documentclass[aspectratio=169, 10pt]{beamer}

% Theme selection
\usetheme{metropolis} % Clean, modern, default light theme

% Packages
\usepackage{listings} % Replaced minted with listings due to local version conflict
\usepackage{graphicx}
\usepackage{booktabs}
\usepackage{hyperref}
\usepackage{amsmath}
\usepackage{xcolor}
\usepackage{tikz}

% Define colors
\definecolor{codebg}{RGB}{245, 245, 245}
\definecolor{keywordcolor}{RGB}{204, 102, 0}
\definecolor{stringcolor}{RGB}{0, 153, 51}
\definecolor{commentcolor}{RGB}{120, 120, 120}

% Configure listings for Python/Text
\lstset{
    backgroundcolor=\color{codebg},
    basicstyle=\ttfamily\small,
    breaklines=true,
    keywordstyle=\color{keywordcolor}\bfseries,
    stringstyle=\color{stringcolor},
    commentstyle=\color{commentcolor}\itshape,
    showstringspaces=false,
    frame=single,
    rulecolor=\color{codebg},
    xleftmargin=1em,
    tabsize=4
}

\title{Demystifying the Lean Theorem Prover}
\subtitle{Reconstructing the Curry-Howard Correspondence via Pure Functional Programming in Python}
\author{Anuj Jha} % Placeholder for student name
\date{\today}
\institute{SLP: Prof. Sudarshan Vijay}

\begin{document}

\maketitle

% -----------------------------------------------------------------------------
\begin{frame}{Agenda}
    \tableofcontents
\end{frame}

% -----------------------------------------------------------------------------
\section{1. The Philosophy: Proofs as Programs}

\begin{frame}{The Core Principle}
    \begin{quote}
        ``Functional programming is one input, one output per function.''
    \end{quote}
    \vspace{1em}
    
    My goal with this project was to understand how modern interactive theorem provers like Lean 4 actually work under the hood. 
    
    What I discovered is that \textbf{Lean's proof system is exactly this concept applied to mathematics.}
    
    \vspace{1em}
    Instead of viewing a proof as a sequence of logical deduction steps, we view a tactic as a pure function:
    \begin{center}
    \texttt{tactic : ProofState $\rightarrow$ ProofState}
    \end{center}
\end{frame}

\begin{frame}{The Curry-Howard Correspondence}
    Why does this work? Because proving a mathematical theorem is literally identical to compiling a computer program.
    
    \vspace{1em}
    \begin{table}
        \centering
        \begin{tabular}{ll}
            \toprule
            \textbf{Mathematics / Logic} & \textbf{Type Theory / Programming} \\
            \midrule
            Proposition (e.g. $2+2=4$) & Type \\
            Proof & A valid value of that type \\
            Tactic & Function that constructs the value \\
            "Q.E.D." (Verification) & Type-Checker accepting the program \\
            \bottomrule
        \end{tabular}
    \end{table}
    \vspace{1em}
    When Lean says "Goals accomplished," it just means the program type-checked successfully.
\end{frame}

% -----------------------------------------------------------------------------
\section{2. Simulating the Lean Kernel in Python}

\begin{frame}[fragile]{Expressions as Tagged Trees}
    In Lean 4, ALL mathematical expressions are represented dynamically as an \texttt{Expr} tree. Since Python lacks algebraic data types, I successfully simulated this using \textbf{immutable tagged tuples}.

\begin{lstlisting}[language=Python]
# Constructors (pure functions that build syntax trees)
def succ(n):   return ("succ", n)                  # S(n)
def add(a, b): return ("add", a, b)                # a + b
def eq(a, b):  return ("eq", a, b)                 # a = b
def le(a, b):  return ("le", a, b)                 # a <= b

# Destructors (For programmatic Pattern Matching)
def is_add(t):   return t[0] == "add"
def left_of(t):  return t[1]
def right_of(t): return t[2]
\end{lstlisting}

    To construct the goal $2 + 2 = 4$:
    \\\texttt{goal = eq(add(two(), two()), four())}
\end{frame}

\begin{frame}[fragile]{The Variable Naming Problem (Concept)}
    A major challenge in math engines is variable names. To a human, \texttt{fun x $\rightarrow$ x} is identical to \texttt{fun y $\rightarrow$ y}. 
    
    But to a computer, comparing the letters 'x' and 'y' means they look completely different!
    
    \vspace{1em}
    \textbf{The Solution: De Bruijn Indices} \\
    Instead of giving variables names, we just give them \textbf{directions}. 
    
    A variable simply becomes: \textit{"I am the variable defined $N$ layers above me."}
\end{frame}

\begin{frame}[fragile]{De Bruijn Indices: A Simple Visual}
    We replace names with "arrows" pointing up to their creator.
    
    \vspace{1em}
    \begin{itemize}
        \item \texttt{fun x $\rightarrow$ x} simply becomes \texttt{fun $\rightarrow$ (Variable 0 layers up)}
        \item \texttt{fun y $\rightarrow$ y} simply becomes \texttt{fun $\rightarrow$ (Variable 0 layers up)}
    \end{itemize}
    
    \vspace{1em}
    Because both functions use exactly the same directions, the computer instantly knows they are mathematically identical! We never have to worry about accidentally mixing up variable names again.
\end{frame}

\begin{frame}[fragile]{Handling Alpha-Equivalence: De Bruijn Indices}
    A major challenge in building a theorem prover is variable shadowing. How does the kernel know \texttt{fun x $\rightarrow$ x} is identical to \texttt{fun y $\rightarrow$ y}?
    
    \vspace{0.5em}
    The solution I studied is the \textbf{De Bruijn Index System}:
    Instead of names, bounded variables (\texttt{BVar}) use numeric offsets pointing upward to the lambda that bounds them.
    
    \vspace{0.5em}
    \begin{itemize}
        \item \texttt{fun x $\rightarrow$ x} internally becomes \texttt{lam("x", BVar(0))}
        \item \texttt{fun y $\rightarrow$ y} internally becomes \texttt{lam("y", BVar(0))}
    \end{itemize}
    
    \vspace{0.5em}
    \textbf{The Result:} Structural hashing entirely ignores the pretty-printing names ("x", "y"). These trees hash to the exact same pointer in memory natively!
\end{frame}

\begin{frame}[fragile]{De Bruijn Indices in Nested Scopes}
    How does this scaling representation work for complex math with multiple variables? The index simply counts how many $\lambda$-binders you must "jump over" going up the AST.
    
    \vspace{1em}
    Consider the standard math expression: \texttt{fun x $\rightarrow$ fun y $\rightarrow$ x (y x)}
    
    \vspace{1em}
    De Bruijn internal representation:
    \begin{center}
    \texttt{lam("x", lam("y", App(BVar(1), App(BVar(0), BVar(1)))))}
    \end{center}
    
    \vspace{1em}
    \begin{itemize}
        \item \texttt{BVar(0)} points to the innermost binder (\texttt{y}).
        \item \texttt{BVar(1)} skips \texttt{y} and points to the outer binder (\texttt{x}).
    \end{itemize}
    This eliminates all alpha-renaming bugs from the Kernel entirely!
\end{frame}

\begin{frame}[fragile]{Rewrite Rules = Pattern-Matching Functions}
    I implemented rewrite rules strictly as purely functional mappings: \texttt{term $\rightarrow$ term | None}. They search for a structural pattern and construct a new AST branch.
    
\begin{lstlisting}[language=Python]
def add_succ_rule(term):
    """ Implementing the Peano Axiom: (a + S(b)) -> S(a + b) """
    if is_add(term) and is_succ(right_of(term)):
        
        a = left_of(term)
        b = pred_of(right_of(term))  
        
        # Return cleanly constructed new tuple
        return succ(add(a, b))       
    return None
\end{lstlisting}

    Notice there is no state mutation. This guarantees mathematical purity.
\end{frame}

\begin{frame}[fragile]{Tactics as Higher-Order Closures}
    To simulate a tactic like Lean's \texttt{rw}, I used \textbf{closures} that take a mathematical rule and return a state transition function.
    
\begin{lstlisting}[language=Python]
def nth_rewrite(n, rule, name):
    """
    Returns a closure that performs DFS traversal, applying the rule 
    to exactly the Nth occurrence it finds.
    """
    def tactic(goal):
        return apply_nth(goal, rule, n)
    tactic._lean = f"nth_rewrite {n} [{name}]"
    
    # Returning the higher-order function composition
    return tactic   
\end{lstlisting}
\end{frame}


% -----------------------------------------------------------------------------
\section{3. Case Study 1: Proving Addition}

\begin{frame}[fragile]{Proving 2 + 2 = 4: The Initial State}
    Let's walk through how my simulation evaluates $2 + 2 = 4$, replicating Lean's behavior exactly.
    
    \vspace{1em}
    \textbf{Initial Goal Tree:} 
    \texttt{eq( add(2, 2), 4 )}
    
    \vspace{1em}
    \textbf{The DFS Trap:} \\
    If we just write \texttt{rw [two\_eq\_succ\_one]}, the DFS traversal hits the VERY FIRST \texttt{2} it finds—the \textbf{left} operand.
    
    Our tree becomes: \texttt{eq( add(S(1), 2), 4 )} $\implies S(1) + 2$
\end{frame}

\begin{frame}[fragile]{Proving 2 + 2 = 4: The Peano Mismatch}
    Why is $S(1) + 2$ a problem?
    
    \vspace{1em}
    Look at the Peano addition axiom (\texttt{add\_succ}):
    \begin{center}
    \Large $a + S(b) \rightarrow S(a + b)$
    \end{center}
    
    \vspace{1em}
    This rule \textbf{ONLY} matches if the \textbf{right} operand has an $S$. It doesn't know what to do with an $S$ on the \textbf{left}! If we unfold the left side, the proof gets completely stuck.
\end{frame}

\begin{frame}[fragile]{Proving 2 + 2 = 4: The Solution}
    We must use \texttt{nth\_rewrite 2}! This forces the engine to bypass the left operand and target the \textbf{second} DFS occurrence.
    
    \vspace{1em}
    \textbf{Step 1:} \texttt{nth\_rewrite(2, two\_eq\_succ\_one)}
    \begin{itemize}
        \item Target the right operand: \texttt{eq( add(2, S(1)), 4 )}
    \end{itemize}
    
    \vspace{1em}
    \textbf{Step 2:} \texttt{rw(add\_succ\_rule)}
    \begin{itemize}
        \item The rule sees the $S$ on the right, and "wraps" it around the addition!
        \item Tree becomes: \texttt{eq( succ(add(2, 1)), 4 )}
    \end{itemize}
    
    This exact peeling process continues 5 more times until we reach \texttt{succ(succ(2)) = 4}.
\end{frame}

% -----------------------------------------------------------------------------
\section{4. Case Study 2: Multiplication \& Decomposition}

\begin{frame}[fragile]{Multiplication is Cascading Addition}
    Multiplication is defined by structural recursion strictly on the right operand:
    
    \begin{equation*}
    a \times 0 = 0
    \end{equation*}
    \begin{equation*}
    a \times \mathrm{succ}(b) = (a \times b) + a
    \end{equation*}
    
    The axiom \texttt{mul\_succ} forces multiplication to decompose into addition subproblems.
    \begin{align*}
        2 \times 2 &= 2 \times S(1) \\
                   &= (2 \times 1) + 2 \\
                   &= ((2 \times 0) + 2) + 2 \\
                   &= (0 + 2) + 2
    \end{align*}
\end{frame}

\begin{frame}[fragile]{Proof Length Asymmetry}
    \textbf{Why does proving $2 \times 2 = 4$ take 19 recursive rewriting steps in my simulation?} \\
    Because proving $(0 + 2) + 2 = 4$ requires cascading applications of \texttt{add\_succ} peeling both \texttt{2} variables apart simultaneously.
    
    \vspace{1em}
    \textbf{A key finding: Commutativity is computationally expensive!}
    Because \texttt{mul\_succ} recurses solely on the \textbf{right}:
    \begin{itemize}
        \item Proving $2 \times 1 = 2$ requires just 10 functional composing steps.
        \item Proving $1 \times 2 = 2$ requires 13 functional composing steps.
    \end{itemize}
    
    Mathematically they are identical. Computationally, their proof structures are entirely distinct.
\end{frame}

% -----------------------------------------------------------------------------
\section{5. Case Study 3: Inequality and Existentials}

\begin{frame}[fragile]{What Actually is an Inequality? (Concept)}
    How do you mathematically prove that $1 \le 2$? 
    
    \vspace{1em}
    In Lean, saying $A \le B$ simply means: \textbf{"There is some GAP number I can add to $A$ to reach $B$."}
    \begin{center}
        \Large $1 \le 2 \implies 1 + \text{gap} = 2$
    \end{center}

    \vspace{1em}
    By defining it this way, if you tell the engine what the gap is, you transform a scary inequality into a simple addition equation!
\end{frame}

\begin{frame}[fragile]{Proving Transitivity: $1 \le 2 \land 2 \le 4 \implies 1 \le 4$}
    Imagine mathematical hypotheses as literal \textbf{packages} handed to you. 
    
    \vspace{1em}
    \textbf{Step 1: Unpacking (\texttt{cases} tactic)} \\
    We use the \texttt{cases} tactic to open our packages.
    \begin{itemize}
        \item Opening the $1 \le 2$ package gives us: a gap of \textbf{1}, and the equation $(2 = 1 + 1)$.
        \item Opening the $2 \le 4$ package gives us: a gap of \textbf{2}, and the equation $(4 = 2 + 2)$.
    \end{itemize}
\end{frame}

\begin{frame}[fragile]{Proving Transitivity: Packing it back up}
    \textbf{Step 2: Packing (\texttt{use} tactic)} \\
    Now we need to prove our final goal: $1 \le 4$.
    
    We just have to provide the total gap! We use the \texttt{use} tactic to say: \textit{"The new gap is just our two old gaps added together!"} ($1 + 2 = \mathbf{3}$).
    
    \vspace{1em}
    Our intimidating $1 \le 4$ goal instantly transforms into:
    \begin{center}
        \Large $4 = 1 + \mathbf{3}$
    \end{center}
    
    We then substitute the equations we unpacked to prove this equation is mathematically true!
\end{frame}

\begin{frame}[fragile]{How is $\le$ ACTUALLY defined?}
    In Lean's Natural Number Game (NNG4), $a \le b$ is not a primitive base truth. It is mathematically defined via an existential type:
    \begin{equation*}
        a \le b \iff \exists c, \; b = a + c
    \end{equation*}

    To prove an inequality, we must provide the "witness" (the numerical gap $c$). By doing so, we essentially translate a complex inequality problem into a simple equality problem that \texttt{rw} can solve!

    In my architecture, the \texttt{use} tactic transforms the inequality goal directly into an equality by injecting this witness.

\begin{lstlisting}[language=Python]
def use(witness, name):
    """ a <= b --> b = a + witness """
    def tactic(goal):
        a = left_of(goal)
        b = right_of(goal)
        return eq(b, add(a, witness))
    return tactic
\end{lstlisting}
\end{frame}

\begin{frame}[fragile]{The Flagship Proof: \texttt{le\_trans} (Modelling Hypotheses)}
    \textbf{Theorem:} \texttt{If 1 <= 2 and 2 <= 4, then 1 <= 4} 
    
    This proof is significantly more advanced because it requires \textbf{Hypotheses}. A theorem is actually just a function taking the assumptions as arguments. Here, the arguments are \texttt{hxy: 1 <= 2} and \texttt{hyz: 2 <= 4}.
    
    I simulated this environment by passing a \texttt{hyps} dictionary payload alongside the goal state throughout the proof timeline:

\begin{lstlisting}[language=Python]
    def proof_le_trans():
        hyps = {
            "hxy": le(one(), two()),    # Granted Hypothesis 1
            "hyz": le(two(), four()),   # Granted Hypothesis 2
        }
        goal = le(one(), four())        # Target Goal
\end{lstlisting}
\end{frame}

\begin{frame}[fragile]{Destructuring Existentials with \texttt{cases}}
    How do we resolve the \texttt{cases} tactic destructuring hypotheses? 
    It executes as a data-manipulating function on the hypothesis dictionary!

    When we call \texttt{cases hxy with a ha}, we are telling the proof state: "Take the fact that $1 \le 2$, give me the gap integer ($a$), and give me the new guaranteed equation ($ha$)."

\begin{lstlisting}[language=Python, basicstyle=\ttfamily\scriptsize]
    # Step 1: cases hxy with a ha
    # Destructs (1<=2) into witness a=1, equation ha: 2=1+1
    def step_cases_hxy(goal, hyps):
        hyps = dict(hyps)
        del hyps["hxy"]  # The old existential hypothesis is consumed
        
        # Inject the inner components of the existential back into scope
        hyps["a"]  = one()
        hyps["ha"] = eq(two(), add(one(), one()))
        return goal, hyps
\end{lstlisting}
\end{frame}

\begin{frame}[fragile]{Bridging the Gap: The \texttt{use} tactic}
    After extracting BOTH witnesses from the hypotheses (we got $a=1$ from $1 \le 2$, and $b=2$ from $2 \le 4$), we must combine them to prove the final goal.
    
    To prove $1 \le 4$, we use the combination $(a + b)$ as our new witness! Since $a=1$ and $b=2$, our combined witness is $3$.

\begin{lstlisting}[language=Python]
    # Step 3: use a + b (combining witnesses)
    def step_use_ab(goal, hyps):
        a, b = hyps["a"], hyps["b"]
        x, z = left_of(goal), right_of(goal)
        
        # Inequality Goal (1 <= 4) structurally transforms into:
        # 4 = 1 + (1 + 2)
        return eq(z, add(x, add(a, b))), hyps  
\end{lstlisting}
\end{frame}

\begin{frame}[fragile]{Hypotheses as Runtime Rewrite Rules}
    Our goal is now \texttt{4 = 1 + (1 + 2)}. How do we verify this? We rewrite the Goal using the actual mathematical data stored inside the hypotheses!

\begin{lstlisting}[language=Python, basicstyle=\ttfamily\scriptsize]
    # Step 4: rw [hb]  
    # Using the hypothesis `hb: 4=2+2` to rewrite 4 in the goal
    def step_rw_hb(goal, hyps):
        # Goal transitions from: 4 = ... to: (2 + 2) = ...
        return eq(add(two(), two()), rhs_of(goal)), hyps

    # Step 5: rw [ha]  
    # Using the hypothesis `ha: 2=1+1` to rewrite the leftmost 2
    def step_rw_ha(goal, hyps):
        # Goal transitions to: ((1 + 1) + 2) = 1 + (1 + 2)
        lhs = lhs_of(goal)
        return eq(add(add(one(), one()), right_of(lhs)), rhs_of(goal)), hyps
        
    # Step 6: exact add_assoc 
    # Terminate by applying the mathematically proven theorem (a+b)+c = a+(b+c)
    def step_add_assoc(goal, hyps):
        return eq(rhs_of(goal), rhs_of(goal)), hyps
\end{lstlisting}
\end{frame}

% -----------------------------------------------------------------------------
\section{6. Insights into the Lean 4 Compiler Architecture}

\begin{frame}[fragile]{The Two-Tier Compiler Pipeline}
    My research into Lean's internals revealed a highly secure separation of concerns:
    
    \begin{center}
    \begin{minipage}{0.85\linewidth}
\begin{lstlisting}[basicstyle=\ttfamily\scriptsize]
     Source Code (.lean)
            |
       [ PARSER ]     Text -> Concrete Syntax Tree
            |
     [ ELABORATOR ]   Syntax -> Core `Expr`. (Where Tactics run!)
            |         Runs `MetaM` Monad, Unification, MetaVariables.
            |         MASSIVE codebase (500,000+ lines). Expected to have bugs.
            |
       [ KERNEL ]     The Trusted Core (~5k lines C++).
            |         Type-checks the generated proof term.
            v
    [ ENVIRONMENT ]   Stores mathematical ground truth permanently.
\end{lstlisting}
    \end{minipage}
    \end{center}
\end{frame}

\begin{frame}{The de Bruijn Criterion \& Complete Isolation}
    Why is the architecture bifurcated this way?
    
    \vspace{1em}
    \begin{itemize}
        \item \textbf{Untrusted Zone:} Tactics, Macros, Parsing, Unification. Allowed to crash, infinite-loop, or have severe logical bugs.
        \item \textbf{Trusted Kernel:} Tiny, mathematically pure type-checker written in memory-safe C++. Analyzes pure ASTs.
        \item If an AI or a buggy tactic generates a fake functional composition, the kernel strictly \textbf{rejects the payload} with a type error. (e.g. \texttt{Array Int (2+2)} failing to reduce against \texttt{Array Int (4)}).
    \end{itemize}
\end{frame}

\begin{frame}{The Type Checker \& Definitional Computation}
    Checking types in Dependent Type Theory is Turing-Complete. \\
    The kernel evaluates expressions via \textbf{Definitional Equality (\texttt{isDefEq})}.
    
    \vspace{1em}
    It features $4$ explicit reduction engine sub-systems:
    \begin{itemize}
        \item \textbf{$\beta$ (Beta) Reduction}: Function execution / lambda evaluation.
        \item \textbf{$\delta$ (Delta) Reduction}: Unfolding definitions (\texttt{double 2 $\rightarrow$ 2 + 2}).
        \item \textbf{$\zeta$ (Zeta) Reduction}: Resolving let-bindings.
        \item \textbf{$\iota$ (Iota) Reduction}: Executing Inductive Type \textbf{Recursors} (Pattern matching).
    \end{itemize}
\end{frame}

% -----------------------------------------------------------------------------
\section{Conclusion}

\begin{frame}{Conclusion: The Ultimate Synthesis}
    \begin{table}
        \centering
        \begin{tabular}{ll}
            \toprule
            \textbf{FP Concept} & \textbf{Lean Architecture Equivalent} \\
            \midrule
            Algebraic Data Types & Tagged Unions (e.g. \texttt{("add", a, b)}) or \texttt{Expr} \\
            Immutability & Modifying proofs returns strictly new nodes \\
            Pattern Matching & Destructing Tuples, executing $\iota$-reduction \\
            Pure Functions & \texttt{goal $\rightarrow$ goal} with zero side effects \\
            Higher-Order Functions & Tactics returning configuration closures \\
            Type Systems & The Curry-Howard Theorem Verification boundary \\
            \bottomrule
        \end{tabular}
    \end{table}
    
    \vspace{1em}
    \begin{center}
        \textbf{Lean operates beautifully on a single mathematical principle: \\ Proving a mathematical theorem IS exactly the compilation of a verified program.}
    \end{center}
\end{frame}

\begin{frame}{Questions?}
    \begin{center}
        \Huge Questions?
        
        \vspace{1em}
        \large Thank you for your time.
    \end{center}
\end{frame}

\end{document}
